La notación asintótica describe cómo crece el costo de un algoritmo cuando el tamaño de la entrada tiende a infinito, pero omite deliberadamente las constantes, la jerarquía de memoria y la estructura particular de los datos. Este informe estudia experimentalmente en qué medida esas omisiones determinan el desempeño observado en dos problemas clásicos: el ordenamiento de un arreglo unidimensional de enteros y la multiplicación de dos matrices cuadradas de enteros \cite{algorithms_erickson}.

Para el problema de ordenamiento se evaluaron MergeSort, QuickSort, PatienceSort y \texttt{std::sort} de la biblioteca estándar de C++; para el de multiplicación, el algoritmo iterativo tradicional y el algoritmo de Strassen. Los cuatro algoritmos de ordenamiento comparten la cota $\mathcal{O}(N \log N)$ en el caso promedio \cite{taocv3}, y los dos algoritmos de multiplicación difieren únicamente en el exponente, $\mathcal{O}(N^{3})$ frente a $\mathcal{O}(N^{2.81})$. Se busca establecer, en cada caso, si el orden de mérito predicho por la teoría se reproduce en la práctica, en qué tamaños de entrada lo hace y qué factores explican las discrepancias. Las mediciones completas que respaldan la discusión se tabulan en el apéndice.
